\documentclass[11pt,letterpaper]{article}
\usepackage[margin=1in]{geometry}
\usepackage{amsmath,amssymb}
\usepackage{hyperref}
\usepackage{enumitem}
\usepackage{parskip}

\title{The Turing: A Unit of Verified Explanatory Gain}
\author{M. Drake Stapleton}
\date{September 29, 2026}

\begin{document}
\maketitle

\noindent\textit{Status: Research hypothesis, v1.4, revised October 2, 2026. Originally published September 29, 2026.}

\noindent\textit{v1.4 revision note: reconciliation pass with the longer papers. Clarifies result status: TY-1+TY-2 is an internally replayed descriptive result; independent blinding is not established.}

\section{The problem}

Artificial intelligence has no unit for understanding gained.

We measure tokens per second. We measure watts. We measure benchmark accuracy. None of these measure how much a machine actually figured out, or what it cost in energy to figure it out.

A model can burn a megawatt to memorize answers. A model can sip power and discover structure that keeps predicting. Today's metrics cannot tell those two apart in any principled way. Accuracy per joule gets close, but accuracy is not understanding: a lookup table can be accurate.

We propose a unit for a measurable component of the missing quantity.

\section{The proposal}

\textbf{The Turing (T)} is a unit of verified explanatory gain.

1 T = 1 bit of net held-out description-length improvement, relative to a declared baseline and a versioned measurement profile, after paying for the explanation itself.

In symbols, for baseline $B$, candidate explanation $M$, sealed held-out observations $D$, declared shared background $G$, and measurement profile $P$:
\begin{equation}
T(M;B,D,G,P) = [L(B|G) + L(D|B,G)] - [L(M|G) + L(D|M,G)]
\end{equation}
the bits to describe each explanation plus the bits to describe the data given it, with the shared background declared rather than smuggled in.

Measurement takes two forms. $T_{\mathrm{ideal}}$ is the complexity-adjusted held-out log-loss improvement: idealized, requiring no actual entropy coder. $T_{\mathrm{real}}$ is the complexity-adjusted improvement in an actual reference-coded bitstream. An idealized result is not a realized compression; the bridge between the two is a claim to be tested, not assumed.

\textbf{Turing yield} is explanatory gain per unit energy: T/J (Turings per joule).

\section{Why the explanation must pay for itself}

A giant lookup table can predict perfectly on the data it stores. Whether that counts for anything depends on what the table had to pay to store it.

Example: a candidate saves 3,000,000 prediction bits but costs 8,000,000 bits to describe. Net: $-5{,}000{,}000$ Turings under profile $P$. It loses.

Another candidate saves 50,000 bits and costs 400 to describe. Net: $+49{,}600$ Turings under profile $P$. It discovered compact, reusable structure.

The second candidate demonstrated greater net held-out explanatory compression under this profile. The mechanism is the complexity charge: it penalizes stored information that does not earn its cost on held-out observations. Under the memorization bound of the formal paper (Lemma~2), pure memorization, explicitly charged, cannot outperform a cheaper equal default on held-out data. That is a deliberately weak claim, and it needs stating carefully: a table-based predictor can legitimately generalize. If a table's entries earn their description cost on observations the table never saw, the table scores well. What is penalized is storage that does not pay its way, not tables as such.

We propose that finding compact, reusable structure which keeps predicting is one measurable component of understanding.

\section{The rules that keep it honest}

\begin{itemize}[leftmargin=*]
\item Held-out data only. Never training data.
\item The baseline is declared up front, not chosen after the fact.
\item Scores are only comparable under the same measurement profile. No context-free leaderboards.
\item Physical energy accounting must conserve: you cannot attribute more joules than were measured.
\item Informational accounting must conserve: you cannot credit more Turings than were produced.
\item Every measurement is replayable from raw evidence by a third party.
\end{itemize}

\section{What this is not}

It is not a replacement for the bit. Bits measure information; Turings measure verified explanatory gain, expressed in bits the way a darwin expresses evolutionary change as a rate (a factor of e per million years).

It is not an intelligence score. A narrow trick can score Turings without being generally intelligent. The claim is deliberately narrow: under profile $P$, candidate $C$ reduced sealed held-out description length by a stated number of bits relative to baseline $B$.

\section{Standing on giants, measuring something new}

The mathematics descends from Shannon's information theory, Solomonoff's algorithmic probability, Kolmogorov complexity, and Rissanen's Minimum Description Length principle (1978), whose two-part code $L(H) + L(D|H)$ is the direct ancestor of the formula above.

None of them named a unit for it. Description length has always been reported in plain bits, which is why a bit of memorization and a bit of genuine explanation look identical on paper.

The contribution is not the name. A name by itself measures nothing. The contribution is the measurement protocol: sealed held-out data, an explanation that pays for itself in the description-length accounting, and a declared profile that makes two reported numbers comparable. There is precedent for naming a unit of this kind: J. B. S. Haldane defined the \textbf{darwin} in 1949 as a unit of evolutionary change, named for Charles Darwin. We define the \textbf{Turing} as a unit of explanatory change.

Related recent work (for example, the Stanford Hazy Research paper ``Intelligence per Watt'' (arXiv:2511.07885), November 2025) measures task accuracy per unit power, and per unit energy in the same work. That is useful and different: accuracy is not explanatory gain, and a memorized answer is accurate. The Turing measures what the accuracy metrics cannot: whether the machine discovered structure that keeps predicting.

\section{Terminology}

\begin{itemize}[leftmargin=*]
\item \textbf{TURING}: the existing AIEN/Omega evidence and decision-record program.
\item \textbf{Turing (T)}: the unit proposed here.
\item \textbf{$T_{\mathrm{ideal}}$}: complexity-adjusted held-out log-loss improvement; idealized, with no entropy coder required.
\item \textbf{$T_{\mathrm{real}}$}: complexity-adjusted improvement in an actual reference-coded bitstream.
\item \textbf{Turing yield}: T/J, explanatory gain per joule.
\end{itemize}

\section{The claim}

To our knowledge, this brief is the first publication of the Turing as a unit of verified explanatory gain, and of Turing yield (T/J) as a measure of the physical efficiency of verified explanatory gain. Published September 29, 2026, by the AIEN project.

We invite everyone to use it. If you report Turings, publish your baseline, your measurement profile, and your evidence, so anyone can recompute your number. That is the entire price of admission.

\section{What comes next}

The first profile-bound idealized result is recorded: TY-1+TY-2, $+2{,}559{,}679.825$ $T_{\mathrm{ideal}}$ under the measurement profile declared for that experiment. That number is an internally reported held-out result under the historical V0 profile, not a fully qualified measurement under the current protocol. The reference-coder bridge remains open: H1, the claim that $T_{\mathrm{ideal}} \approx T_{\mathrm{real}}$ within the declared bound, is not yet verified. Results will be published as they verify, including failures.

\bigskip
\noindent\textit{One measurable component of understanding is the discovery of compact structure that keeps predicting what happens next. Now it has a unit.}

\end{document}
